% =============================================================================
\section{Regional Allocation Model}\label{sec:optimisation}
% =============================================================================

The allocation module decides how many surge beds to place in each NHS
region and how to route excess demand through inter-regional patient
transfers, given the quantile forecasts of \Cref{sec:forecasting}. We
present a deterministic linear program, a risk-averse (tail-weighted) extension,
and closed-form heuristic baselines.

\subsection{Sets, Parameters, and Decision Variables}\label{ssec:notation}

The scope is \emph{hybrid regional}, with demand nodes and capacity sites
coinciding with the seven NHS England commissioning regions, one set
$\R=\{r_1,\ldots,r_7\}$. Unlike facility-location models that choose
which sites to open~\cite{bertsimas2022where}, every region already
operates critical care, so there is no opening decision and surge cost
is uniform per bed; the model decides only bed counts and routing. The indices are regions $r,r'\!\in\!\R$; horizons
$h\!\in\!\Hset=\{7,14,21,28\}$~days; scenarios
$s\!\in\!\Sset=\{\text{low},\text{median},\text{high}\}$ from
$q^{(0.1)},q^{(0.5)},q^{(0.9)}$ with weights
$\pi=(0.2,0.6,0.2)$, a representative-quantile discretisation of lower,
central, and upper demand stress (a scenario mixture, not calibrated
probability bins).
The parameters are forecast demand
$\xi_{r,h}^s=\hat q_r^{(s)}(t_0+h)$; baseline critical-care capacity
$\Gamma_r=1.05\times$Delta-peak;
per-region surge cap
$K_r=0.5\,\Gamma_r$; budget
$B=0.20\sum_r \Gamma_r$; great-circle inter-centroid distance $\ell_{r r'}$ in
km ($\times 1.3$ road-detour); travel time
$T_{r r'}=60\,\ell_{r r'}/80$~min; mutual-aid cap $\tau=240$~min; outbound transfer cap $\kappa$
(uncapped at the headline operating point).
Decision variables fall in two layers. The \emph{first-stage} capacity
decision, the planner's only true control, is the surge allocation
$b_{r}\!\in\![0,K_r]$, shared across scenarios and horizons. The
per-scenario \emph{second-stage} recourse is the flow
$z_{r r',h}^s\!\geq\!0$ (defined for $T_{r r'}\!\leq\!\tau$); its
self-loop $z_{r r,h}^s$ is demand served \emph{locally} and its
off-diagonal ($r'\!\neq\!r$) entries are inter-regional transfers, with
unmet demand $u_{r,h}^s\!\geq\!0$ and a tail auxiliary $W\!\geq\!0$ for
the risk-averse variant. Of the $42$ ordered region pairs, $24$ satisfy
$T_{r r'}\!\leq\!\tau$ ($12$ undirected mutual-aid links); Midlands
reaches every region and South West only Midlands.

\subsection{Deterministic Median-Scenario LP}\label{ssec:milp}

For the median scenario $m$ ($q^{(0.5)}$), the deterministic model solves
\begin{equation}\label{eq:deterministic_objective}
\min\;
\textstyle\sum_{r,h} u_{r,h}^{m}
\;+\;\theta\sum_{r\neq r',h} \ell_{r r'}\, z_{r r',h}^{m},
\end{equation}
with $\theta=10^{-3}$ a transfer tie-break weight (one bed of weekly
snapshot unmet outweighs $1000$~patient-km of transfer). The objective
minimises \emph{unmet demand}; the second term only breaks ties between
equal-coverage plans by preferring less patient transport, so the surge
beds $b$ are placed to reduce shortage, not to minimise transfers. Each
$u_{r,h}^{m}$ is a snapshot of bed-shortfall at the weekly checkpoint
$h$, so the objective is the snapshot unmet summed over the four-week
planning horizon, not a continuous bed-day integral. The model is subject to
\begin{align}
z_{r r,h}^{m}+\!\!\textstyle\sum_{r'\neq r:\,T_{r r'}\leq\tau}\! z_{r r',h}^{m} + u_{r,h}^{m}
  &= \xi_{r,h}^{m}
  && \forall r,h, \label{eq:demand}\\
z_{r r,h}^{m}+\!\!\textstyle\sum_{r'\neq r:\,T_{r' r}\leq\tau}\! z_{r' r,h}^{m}
  &\leq \Gamma_r + b_{r}
  && \forall r,h, \label{eq:capacity}\\
\textstyle\sum_{r'\neq r}\, z_{r r',h}^{m} &\leq \kappa
  && \forall r,h, \label{eq:transfercap}\\
\textstyle\sum_r b_r &\leq B,
  && \label{eq:budget}\\
z_{r r',h}^{m} &= 0
  && \text{if } T_{r r'}\!>\!\tau, \label{eq:travel}\\
b_{r},z_{r r',h}^{m},u_{r,h}^{m} &\geq 0.
\end{align}
\Cref{eq:demand} is a demand balance. Each region's forecast demand is
served locally ($z_{r r,h}^{m}$), transferred to a reachable region
($z_{r r',h}^{m}$, $r'\!\neq\!r$), or left unmet ($u_{r,h}^{m}$).
\Cref{eq:capacity} caps the beds a region uses (local plus
transferred-in patients) at baseline plus surge; \Cref{eq:transfercap}
bounds outbound patients by the transfer-service capacity $\kappa$;
\Cref{eq:budget} charges the budget once per region; \Cref{eq:travel}
suppresses transfers beyond the mutual-aid window. As $u_{r,h}^{m}$ is
minimised, \Cref{eq:demand} holds it at the residual shortfall
$\max(0,\, \xi_{r,h}^{m}-\sum_{r'} z_{r r',h}^{m})$, a derived quantity
rather than a free decision.

\subsection{Risk-Averse Tail-Weighted Objective}\label{ssec:robust}

To hedge against the high-demand scenario the risk-averse model imposes
\Cref{eq:demand,eq:capacity,eq:transfercap} for every $s\in\Sset$,
retains \Cref{eq:budget,eq:travel}, and adds a tail term $\lambda_3 W$
($W\geq\sum_{r,h} u_{r,h}^s$ for $s\in\Sset_{\text{tail}}=\{\text{high}\}$),
giving the objective
\begin{equation}\label{eq:robust_objective}
\min\;
\sum_{s\in\Sset}\!\pi_s\!\Big(
\textstyle\sum_{r,h} u_{r,h}^s + \theta \sum_{r\neq r',h} \ell_{r r'} z_{r r',h}^s
\Big)\;+\;\lambda_3\, W.
\end{equation}
This is a \emph{scenario-weighted} LP that up-weights the high-demand
tail, not robust optimisation in the formal sense. With a single tail
scenario, minimisation drives $W=\sum_{r,h}u_{r,h}^{\text{high}}$, so
$\lambda_3 W$ is a linear penalty on the worst-scenario shortfall
(equivalently, extra weight on the high scenario) rather than a
Rockafellar--Uryasev conditional value-at-risk (CVaR); the auxiliary
form generalises unchanged to multiple tail scenarios. We call it the
\emph{robust LP} for brevity and report $\lambda_3{=}1$ throughout
($\lambda_3{=}0$ recovers the risk-neutral objective and
$\lambda_3\!\to\!\infty$ a worst-case policy on the high scenario); the
tail term binds only under tight budgets (\Cref{ssec:e5_e6}). Imposing
$q^{(0.9)}$ jointly across regions and horizons makes the construction
co-monotone, its tail over-conservative relative to the true joint
distribution; we treat this as an explicit stress-test choice and
examine its efficiency cost in \Cref{ssec:e5_e6}.

\subsection{Exact LP Solution and Baselines}\label{ssec:solvers}

Both models are linear programs ($7$ surge variables, $84$ self-loops,
up to $504$ transfer flows, $84$ unmet variables, and the tail
auxiliary, all continuous); a deployed plan rounds the per-region surge
to integers, changing total coverage by at most a few beds at this
budget. We solve them exactly in PuLP with the bundled COIN-OR
Branch-and-Cut (CBC) solver; at seven-region scale each instance reaches
global optimality in $<\!0.1$~s, so metaheuristic search is deferred to
the trust-level extension (\Cref{sec:conclusion}).

Every policy is scored on equal footing by a common LP re-evaluation.
For a fixed surge $\mathbf{b}^\star$ ($\sum_r b^\star_r\leq B$,
$b^\star_r\leq K_r$), the residual problem in $(z,u)$ fills in transfers
and unmet demand across all three scenarios, giving honest $E[u]$ and
$u^{\text{worst}}$ for the deterministic, robust, and heuristic
allocations alike.

Finally, three closed-form heuristics give a no-optimisation floor.
Population-proportional uses ONS region populations,
demand-proportional uses expected peak demand, and greedy
shortage-first adds single-bed increments to the most-unmet
region. Each fixes $\mathbf{b}^\star$, completed by the same LP step.
